\begin{lemma}
Let \(K\subseteq \operatorname{Fin}(m)\times \operatorname{Fin}(m)\) have
cardinality \(m\), and suppose \(n\le m^2\).  For every threshold \(T>0\),
the number of injections \(\pi:\operatorname{Fin}(n)\hookrightarrow
\operatorname{Fin}(m)\times\operatorname{Fin}(m)\) for which the sampled image
\(\pi(\operatorname{Fin}(n))\) meets \(K\) in more than \(T\) points is at most
\[
\left(\frac{e(n/m)}{T}\right)^T
\cdot
\#\{\pi:\operatorname{Fin}(n)\hookrightarrow
\operatorname{Fin}(m)\times\operatorname{Fin}(m)\}.
\]
\end{lemma}
\begin{proof}
Apply \(\noderef{UniformInjectionImage}\) to the full domain
\(\operatorname{Fin}(n)\) and to the predicate that an \(n\)-element image set
meets \(K\) in more than \(T\) points.  This identifies the injection count,
up to the common size of each image fiber, with the corresponding count of
\(n\)-subsets of \(\operatorname{Fin}(m)\times\operatorname{Fin}(m)\).
Transport those subsets across the standard equivalence with
\(\operatorname{Fin}(m^2)\), using \(\noderef{FinProductTailCountEquiv}\).
The image of \(K\) has cardinality \(m\), so the resulting subset count is a
hypergeometric upper tail with population \(m^2\), sample size \(n\), marked
set size \(m\), and mean \(n/m\).  The desired bound is exactly
\(\noderef{HypergeometricCountingTailBound}\), and multiplying back by the
number of injections gives the displayed estimate.
\end{proof}
